\documentclass[11pt]{article}

\usepackage[a4paper,margin=1in]{geometry}
\usepackage{amsmath,amssymb,amsthm,mathtools}
\usepackage[hidelinks]{hyperref}
\hypersetup{
  pdfauthor={Bangyen Pham},
  pdftitle={Coefficient Mass of Polynomial Multiples at Complex Roots},
  pdfsubject={Lower bounds for the coefficient order statistics and logarithmic coefficient mass of polynomial multiples with prescribed complex roots},
  pdfkeywords={polynomial multiples, coefficient mass, complex roots, Newton polygon, Gaussian integers},
  pdflang={en-US}
}

\newtheorem{theorem}{Theorem}[section]
\newtheorem{lemma}[theorem]{Lemma}
\newtheorem{corollary}[theorem]{Corollary}
\newtheorem{proposition}[theorem]{Proposition}
\newcommand{\R}{\mathbb R}
\newcommand{\N}{\mathbb N}

\title{Coefficient Mass of Polynomial Multiples at Complex Roots}
\author{\shortstack{Bangyen Pham\\
\small Independent Researcher\\
\small\href{mailto:bangyenp@gmail.com}{\texttt{bangyenp@gmail.com}}\\
\small\href{https://orcid.org/0009-0006-9928-2088}{ORCID: 0009-0006-9928-2088}}}
\date{September 2026}

\begin{document}
\maketitle

\begin{abstract}
Let $b_k(F)$ be the $k$-th largest nonleading coefficient magnitude of
$F=\sum_{j\le D}f_jx^j$ and $\Lambda(F)=\sum_j\log^+|f_j|$.  At real roots
$r_i\ge2$ a companion paper gives $b_k(F)\ge|f_D|\prod_{i\ge k}(r_i-1)$ for
every multiple and charges the $i$-th smallest root $i$ times in $\Lambda$.
Purely imaginary pairs $\pm ic$ reduce by multisection exactly to real roots
$c^2$: every $b_k$ and $\Lambda$ have equal infima over monic multiples.
Otherwise only the first row survives: $b_2$ can vanish and the mass be
linear for pairs within any fixed angle of the imaginary axis.  An
Archimedean Newton polygon restores rows and mass up to constants at moduli
separated by a factor $9$; with several roots per annulus, separation $162$
gives a sharp mass bound charging each root once per annulus at or below
it, and the real-root cross terms fail.  Real roots of both signs charge one root
of each sign per excluded coefficient, with sharp mass constant $\frac14$.
At distinct Gaussian integers of modulus at least $\rho>2$ above the
axis, integer multiples split into blocks paying $\log(\rho/2)$ per unit of
degree.
\end{abstract}

\noindent\textbf{Keywords.} Polynomial multiples; coefficient mass; complex
roots; Newton polygon; Gaussian integers.

\noindent\textbf{2020 Mathematics Subject Classification.} Primary 11C08; Secondary 30C15, 26C10.

\section{Introduction}

For a nonzero $F(x)=\sum_jf_jx^j$ of degree $D$ write
\[
  \Lambda(F)=\sum_{j:f_j\ne0}\log^+|f_j|,\qquad \log^+x=\max(0,\log x),
\]
for its logarithmic coefficient mass, and $b_k(F)$ for the $k$-th largest of
the $D$ nonleading magnitudes $|f_j|$, $j<D$, counted with repetition and
zeros included ($b_k(F)=0$ for $k>D$).

For real roots
$2\le r_1\le\cdots\le r_L$, repetitions allowed, the companion
paper~\cite{coefficient-mass} proves that every nonzero real or complex
multiple $F$ of $\prod_{i=1}^L(x-r_i)$ satisfies
\begin{equation}\label{eq:order}
  b_k(F)\ \ge\ |f_D|\prod_{i=k}^L(r_i-1)\quad(1\le k\le L),
  \qquad
  \Lambda(F)\ \ge\ (L+1)\log|f_D|+\sum_{i=1}^Li\log(r_i-1)
\end{equation}
\cite[Theorem~3.2, Corollaries~3.3 and~3.4]{coefficient-mass}, uniformly in
the degree; we call the first bound at a given $k$ the row $k$.  The row
$k=1$ needs only $r_1>1$.  This paper asks what remains of
\eqref{eq:order} when the prescribed roots are complex.

\paragraph{Main results.}  Purely imaginary pairs transfer exactly: a pair
$\pm ic$ is worth one real root $c^2$, with the same infima of every $b_k$
and of $\Lambda$ over monic multiples (Theorem~\ref{thm:transfer},
Corollary~\ref{cor:imag}).  At arbitrary angles only the first row survives:
it holds at every root set, while $b_2$ can vanish and the mass be linear in
the number of pairs (Propositions~\ref{prop:rowone} and~\ref{prop:noangle}).
Moduli separated by a factor $9$ restore the rows and the mass up to
constants (Corollary~\ref{cor:separated}); with several roots per annulus,
separation $162$ gives a mass bound, sharp to leading order, that charges
each root once for every annulus at or below it, where the cross terms of
\eqref{eq:order} fail (Theorem~\ref{thm:fixedgap},
Proposition~\ref{prop:annulisharp}).  Real roots of both signs give a mass
bound with the sharp leading constant $\frac14$
(Corollary~\ref{cor:bothsignsmass}, Proposition~\ref{prop:bothsignssharp}).
At distinct Gaussian-integer roots of modulus at least $\rho>2$ above the
axis, integer multiples split into blocks paying $\log(\rho/2)$ per unit of
degree (Theorem~\ref{thm:gausscharge}).

Reduced height and reduced length, the least height and length of a
normalized multiple, have been studied for real and complex
coefficients~\cite{dj,schinzel-complex}, and power series with restricted
coefficients and a root on a given ray in~\cite{bbbp-ray}.  The first row
below, $b_1(F)\ge|f_D|\prod_j(|\alpha_j|-1)$, is the complex analogue of the
bound $H(P)\ge|P(1)|$ on the reduced height $H(P)$ (the infimum of the
heights of the multiples $PQ$, $Q$ monic) obtained in~\cite{dj} for real roots
greater than one under an extra sign condition~\cite[p.~333]{dj}; that
condition is removed by the $u=0$ clause of~\cite[Theorem~3.2]{coefficient-mass}.
The two bounds agree at such roots, but at complex
roots $|P(1)|$ is not a lower bound for $H(P)$: $P=x^2+\rho^2$ has $H(P)\le\rho^2$ (take $Q=1$),
while $|P(1)|=1+\rho^2$ and $\prod_j(|\alpha_j|-1)=(\rho-1)^2$.

\paragraph{Exact reductions.}  If the prescribed factor has the form
$R(x^m)$, multisection reduces every multiple to a multiple of $R$ with a
subset of the coefficients (Theorem~\ref{thm:transfer}).  In particular a
purely imaginary pair $\pm ic$ is worth exactly one real root $c^2$, and all
the real-root results transfer (Corollary~\ref{cor:imag}).  A rotation
reduces roots on a common ray to real roots at their moduli, which gives the
rows (Corollary~\ref{cor:ray}).

\paragraph{Arbitrary angles.}  The first row survives at every root set
(Proposition~\ref{prop:rowone}), but nothing more does in general: pairs
within any fixed angle of the imaginary axis admit multiples with $b_2=0$
and mass linear in the number of pairs (Proposition~\ref{prop:noangle}).

\paragraph{Separated moduli.}  Moduli separated by a factor $9$ restore the
rows and the mass up to constants (Corollary~\ref{cor:separated}), and with
several roots per annulus a separation $162$ between annuli gives a mass
bound charging each root once for every annulus at or below it
(Theorem~\ref{thm:fixedgap}).  Both come from an Archimedean Newton polygon
(Theorem~\ref{thm:archnewton}).  In detail, for $K_s$ roots with moduli in $[T_s,U_s]$,
$1\le s\le S$, where $T_1>3$ and $T_s>162U_{s-1}$, the central index charges a
root once for every annulus at or below it,
$\sum_ssK_s\log(T_s/2)-O(S)$ in all for multiples with $|f_D|\ge1$, attained for suitable roots as $\min_sT_s\to\infty$; the cross terms $K_uK_s\log T_s$ of \eqref{eq:order} fail
(Theorem~\ref{thm:annuli},
Corollary~\ref{cor:annuli}, Theorem~\ref{thm:fixedgap},
Proposition~\ref{prop:annulisharp}).  A central index at radius $r$ has at
least $n$ positions above it once $n$ roots have modulus at least $3nr$
(Lemma~\ref{lem:tropcount}).

\paragraph{Both signs.}  For real roots of both signs, Rolle's theorem on
each half-line charges one root of each sign per excluded coefficient
(Theorem~\ref{thm:bothsigns}); the resulting mass bound has the sharp
leading constant $\frac14$ (Corollary~\ref{cor:bothsignsmass},
Proposition~\ref{prop:bothsignssharp}).

\paragraph{Gaussian roots.}  At distinct Gaussian-integer roots in the upper
half-plane, an integer multiple whose coefficients below a position are
small has its part from that position on vanishing at the roots as well
(Lemma~\ref{lem:gausscarry}).  So it splits into blocks
(Lemma~\ref{lem:gaussblocks}), each paying $\log(\rho_{\min}/2)$ per unit of
degree, whatever the degree (Theorem~\ref{thm:gausscharge}).  Roots in a
sector, where the argument principle along a ray next to the sector gives
many real zeros of a real polynomial dominated coefficientwise by the
multiple, are treated in~\cite{coefficient-mass-sectors}, which uses
Theorem~\ref{thm:gausscharge} for integer multiples of large degree.

\paragraph{Open questions.}  Two questions are left open.  First, whether
the rows of Corollary~\ref{cor:annuli}, item~1, hold at a fixed separation
of the annuli, and what the least fixed separation for the mass is, in
particular whether the separation $9$ of Corollary~\ref{cor:annuli}, item~2,
suffices (Section~\ref{sec:annuli}).  Second, whether
$\Lambda(F)\ge cK^2\log R$ holds for every integer multiple $F$ at $K$
distinct Gaussian-integer roots of modulus at least $R$ in the upper
half-plane within a sector of fixed half-width $\delta<\pi/2$.  For real multiples it
fails already at $x^{2M}+\rho^{2M}$ (Proposition~\ref{prop:noangle}); for
integer multiples at Gaussian roots in such a sector the bound
of~\cite{coefficient-mass-sectors} is $(\frac\pi\delta-o(1))K\log R$, linear
in $K$.

\paragraph{Conventions.}
All logarithms are natural.  For $F\in\mathbb C[x]$ of degree $D$ we write
$F^*(t)=t^DF(1/t)=\sum_{d=0}^Df_{D-d}t^d$ for the reciprocal polynomial.
Coefficients of a polynomial named by a capital letter, when used, carry the
lowercase letter: $F=\sum_jf_jx^j$, $B=\sum_ib_ix^i$.  With an argument,
$b_k(\cdot)$ is always the $k$-th largest nonleading magnitude, so $b_k(B)$
and the coefficient $b_k$ of $B$ differ.
References of the form \cite[Theorem~3.2]{coefficient-mass} are to the
companion paper.

\section{Exact reductions}\label{sec:exact}

\begin{lemma}[Multisection]\label{lem:multisection}
Let $m\ge1$, let $R\in\mathbb C[y]$ with $R(0)\ne0$, and put $P(x)=R(x^m)$.  Write
a polynomial $F\in\mathbb C[x]$ as
\[
  F(x)=\sum_{\varepsilon=0}^{m-1}x^\varepsilon E_\varepsilon(x^m),
  \qquad E_\varepsilon(y)=\sum_{i\ge0}f_{mi+\varepsilon}\,y^i .
\]
If $P$ divides $F$, then $R$ divides every $E_\varepsilon$.  If $F$ and $R$
are real, so are the $E_\varepsilon$.
\end{lemma}

\begin{proof}
Let $\zeta=e^{2\pi i/m}$.  Since $P(\zeta^kx)=R(\zeta^{km}x^m)=P(x)$, the
polynomial $P$ divides $F(\zeta^kx)$ for every $k$, hence divides
\[
  \frac1m\sum_{k=0}^{m-1}\zeta^{-k\varepsilon}F(\zeta^kx)
  =x^\varepsilon E_\varepsilon(x^m).
\]
As $R(0)\ne0$, $P(0)\ne0$, so $x^\varepsilon$ is coprime to $P$ and
$R(x^m)\mid E_\varepsilon(x^m)$.  Divide in $\mathbb C[y]$:
$E_\varepsilon=RQ+S$ with $\deg S<\deg R$.  Then $R(x^m)$ divides
$S(x^m)$, whose degree $m\deg S$ is smaller than $m\deg R$, so $S=0$.
Reality of $E_\varepsilon$ is clear from its coefficients.
\end{proof}

\begin{theorem}[Multisection transfer]\label{thm:transfer}
Let $m$, $R$ and $P=R(x^m)$ be as in Lemma~\ref{lem:multisection}, and let
$F$ be a nonzero multiple of $P$ of degree $D$.  For the residue
$\varepsilon\equiv D\pmod m$, the polynomial $E_\varepsilon$ is a multiple of
$R$ of degree $(D-\varepsilon)/m$ with leading coefficient $f_D$, and
\[
  b_k(F)\ge b_k(E_\varepsilon)\quad(k\ge1),\qquad
  \Lambda(F)\ge\Lambda(E_\varepsilon).
\]
Conversely, for every multiple $G$ of $R$ the polynomial $G(x^m)$ is a
multiple of $P$ with the same nonzero coefficients.  Consequently, for
$\Phi=b_k$ (any fixed $k$) and for $\Phi=\Lambda$,
\[
  \inf\{\Phi(F):F\text{ monic},\ P\mid F\}
  =\inf\{\Phi(G):G\text{ monic},\ R\mid G\},
\]
over complex multiples, and also over real multiples when $R$ is real.
\end{theorem}

\begin{proof}
By Lemma~\ref{lem:multisection}, $R\mid E_\varepsilon$.  The top coefficient
of $F$ sits at $D=m\cdot\frac{D-\varepsilon}m+\varepsilon$, so it is the
coefficient of $y^{(D-\varepsilon)/m}$ in $E_\varepsilon$, and no higher
coefficient of $E_\varepsilon$ is nonzero.  The nonleading coefficients of
$E_\varepsilon$ are the $f_j$ with $j\equiv\varepsilon$, $j<D$: a
sub-multiset of the nonleading coefficients of $F$.  The $k$-th largest
element of a sub-multiset is at most the $k$-th largest of the whole, and
$\Lambda$ is a sum of nonnegative terms over coefficients; this gives the
two inequalities.  Monicity passes from $F$ to $E_\varepsilon$, which gives
$\ge$ between the infima.  For the converse, $R\mid G$ gives
$R(x^m)\mid G(x^m)$, and $G(x^m)$ has the coefficients of $G$ spread out
with zeros between them, so $b_k$ and $\Lambda$ are unchanged, and
$G(x^m)$ is monic when $G$ is.  This gives $\le$.
\end{proof}

The case $m=2$, $R(y)=\prod_j(y+c_j^2)$ is the case of purely imaginary
conjugate pairs; the sign change $y\mapsto-y$ turns $R$ into a product of
linear factors at the positive roots $c_j^2$ without changing any
coefficient magnitude.

\begin{corollary}[Purely imaginary pairs]\label{cor:imag}
Let $c_1\le\cdots\le c_K$ be positive reals, repetitions allowed, put
$P_{\rm im}=\prod_{j=1}^K(x^2+c_j^2)$ and
$P_{\rm re}=\prod_{j=1}^K(y-c_j^2)$, and let $F\in\mathbb C[x]$ be a nonzero
multiple of $P_{\rm im}$.
\begin{enumerate}
\item If $c_1^2\ge2$, then for $1\le k\le K$,
\[
  b_k(F)\ \ge\ |f_D|\prod_{j=k}^K(c_j^2-1),
\]
and at least $k$ positions $j<D$ with $j\equiv D\pmod2$ carry magnitudes
$|f_j|$ at least this bound.
The row $k=1$ holds for $c_1>1$.
\item If $c_1^2\ge2$ and $|f_D|\ge1$, then
\[
  \Lambda(F)\ \ge\ (K+1)\log|f_D|+\sum_{j=1}^Kj\log(c_j^2-1).
\]
\item For $\Phi=b_k$ and $\Phi=\Lambda$, the infimum of $\Phi$ over monic
(real or complex) multiples of $P_{\rm im}$ equals its infimum over monic
multiples of $P_{\rm re}$.  In particular, with $c_1^2\ge2$ and $K\ge2$,
\[
  \inf_{Q\text{ monic}}\Lambda(P_{\rm im}Q)\ \asymp\
  K^2+\sum_{j=1}^Kj\log(c_j^2-1)
\]
with the absolute constants of \cite[Corollary~4.3]{coefficient-mass};
$\Lambda(F)\ge K^2/80$ for monic $F$ when $K\ge8$; and for $K=2$ and
$(c_1^2,c_2^2)=(2,3)$ the row $k=2$ of item~1 is sharp:
$\inf_{Q\text{ monic}}b_2(P_{\rm im}Q)=c_2^2-1=2$.
\end{enumerate}
\end{corollary}

\begin{proof}
Apply Theorem~\ref{thm:transfer} with $m=2$ and $R(y)=\prod_j(y+c_j^2)$;
$R(0)\ne0$.  The multiple $E_\varepsilon$ of $R$ has leading coefficient
$f_D$, and $\tilde E(y)=E_\varepsilon(-y)$ is a multiple of $P_{\rm re}$
with leading coefficient $\pm f_D$ and the same coefficient magnitudes.  \cite[Theorem~3.2]{coefficient-mass}, with \cite[Corollary~3.3]{coefficient-mass} for complex $F$, applied to
$\tilde E$ at the roots $c_j^2\ge2$ gives item~1, since
$b_k(F)\ge b_k(\tilde E)$ and the positions of $\tilde E$ are those
$j\equiv D\pmod 2$ of $F$; the clause $c_1>1$ is the clause $r_1>1$ for the row $k=1$ of
\cite[Theorem~3.2]{coefficient-mass}.  Item~2 is \cite[Corollary~3.4]{coefficient-mass} applied to $\tilde E$, together with
$\Lambda(F)\ge\Lambda(\tilde E)$.  Item~3 is the equality of infima in
Theorem~\ref{thm:transfer}, composed with the magnitude-preserving
bijection $G(y)\mapsto\pm G(-y)$ between monic multiples of $R$ and of
$P_{\rm re}$; the three quoted statements are \cite[Corollary~4.3]{coefficient-mass},
\cite[Proposition~4.2]{coefficient-mass} and \cite[Proposition~4.4]{coefficient-mass} at
the roots $c_j^2$.  These are stated for real multiples; they hold for complex
ones because, as in \cite[Corollary~3.3]{coefficient-mass}, a monic complex
multiple $G$ of $P_{\rm re}$ dominates coefficientwise the monic real multiple
$\operatorname{Re}G$, and $b_k$ and $\Lambda$ do not increase under such
domination.
\end{proof}

Thus a purely imaginary pair $\pm ic$ is worth exactly one real root $c^2$:
the pair has two roots of modulus $c$, and the mass charges $\log c^2$ once
at its rank.  This is the right count.  By item~3 the product
$P_{\rm im}$ itself is within $O(K^2)$ of the infimum
(\cite[Corollary~3.7]{coefficient-mass}), and $\Lambda(P_{\rm im})=\Lambda(P_{\rm re})$; so a
bound of the form $\sum_{i\le2K}i\log\rho_{(i)}$, where
$\rho_{(1)}\le\cdots\le\rho_{(2K)}$ are the moduli of the $2K$ roots
counted separately, which would be about twice as large, is false for
$P_{\rm im}$ once $\min_jc_j$ is large: it exceeds the bound
$\Lambda(P_{\rm im})\le\sum_j\log\binom Kj+\sum_j2j\log c_j$ by at least
$K^2\log c_1-\sum_j\log\binom Kj$.

The same proof applies to every factor of the form $R(x^m)$ with $R$ real
and all roots of $R$ of one sign: the real roots $\pm r$ in pairs
($R=\prod(y-r_j^2)$), or, for fixed $m$, all $m$-th roots of real $s_j$ of
one sign.  One
more reduction needs no multisection.

\begin{corollary}[Roots on a common ray]\label{cor:ray}
Let $\theta\in\R$ and $\rho_1\le\cdots\le\rho_K$ with $\rho_1\ge2$, and let
$F\in\mathbb C[x]$ be a nonzero multiple of $\prod_j(x-\rho_je^{i\theta})$.  Then
$b_k(F)\ge|f_D|\prod_{j\ge k}(\rho_j-1)$ for $1\le k\le K$.  In
particular this holds for a real $F$ divisible by
$\prod_j(x^2-2\rho_j\cos\theta\,x+\rho_j^2)$.
\end{corollary}

\begin{proof}
$G(x)=F(e^{i\theta}x)$ has coefficients $f_je^{ij\theta}$, of the same
magnitudes, and is divisible by $\prod_j(x-\rho_j)$; apply
\cite[Corollary~3.3]{coefficient-mass}.
\end{proof}

A common ray charges each pair once, with $\log(\rho-1)$ in place of
$\log(\rho^2-1)$; the imaginary axis is special because it carries both
members of each pair.


\section{Arbitrary angles}\label{sec:angles}

\begin{proposition}[The first row at arbitrary roots]\label{prop:rowone}
Let $\alpha_1,\ldots,\alpha_N\in\mathbb C$, repetitions allowed, with
$\rho_j=|\alpha_j|$, and let $F\in\mathbb C[x]$ be a nonzero multiple of
$\prod_j(x-\alpha_j)$.
\begin{enumerate}
\item If every $\rho_j>1$, then $b_1(F)\ge|f_D|\prod_{j=1}^N(\rho_j-1)$.
\item For every $0<r<1$,
$b_1(F)\ge|f_D|\,\frac{1-r}{r}\bigl(r^N\mathcal M-1\bigr)$, where
$\mathcal M=\prod_j\max(1,\rho_j)$; with $r=N/(N+1)$,
$b_1(F)\ge|f_D|\,(\mathcal M/e-1)/N$.
\end{enumerate}
\end{proposition}

\begin{proof}
Let $F^*(t)=t^DF(1/t)=\sum_{d=0}^Df_{D-d}t^d$ and $\beta_j=1/\alpha_j$
where $\alpha_j\ne0$.

\emph{Item 1.}  Here every $\alpha_j\ne0$, and if $(x-\alpha)^\mu\mid F$
then $F^*(t)=(1-\alpha t)^\mu\,t^{D-\mu}H(1/t)$ with $F=(x-\alpha)^\mu H$, so
$\chi(t)=\prod_j(t-\beta_j)$ divides $F^*$.  Put
$B(s)=\prod_j(1-\beta_js)=s^N\chi(1/s)$, let $c=-[s^N]B$, and let
$v_d=[s^d]V$ for $V=(B+cs^N)/B=1+cs^N/B$.  The numerator $A=B+cs^N$ has
degree at most $N-1$, so $BV=A$ gives, for every $n\ge N$,
$\sum_{k=0}^N[s^k]B\;v_{n-k}=0$, that is,
$\sum_{i=0}^N\chi_iv_{m'+i}=0$ for every $m'\ge0$, where
$\chi=\sum_i\chi_it^i$.  Writing $F^*=\chi h$,
\[
  \sum_{d}v_df_{D-d}
  =\sum_{m'}h_{m'}\sum_i\chi_iv_{m'+i}=0 .
\]
Since $v_0=1$, $|f_D|\le b_1(F)\sum_{d\ge1}|v_d|$.  Now
$|c|=\prod_j|\beta_j|=\prod_j\rho_j^{-1}$, and the coefficients of $1/B$ are
dominated in modulus by those of $\prod_j(1-|\beta_j|s)^{-1}$, whose sum is
$\prod_j(1-\rho_j^{-1})^{-1}$.  Hence
$\sum_{d\ge1}|v_d|\le\prod_j\rho_j^{-1}(1-\rho_j^{-1})^{-1}
=\prod_j(\rho_j-1)^{-1}$.

\emph{Item 2.}  $F^*(0)=f_D\ne0$, and $F^*$ vanishes at every $\beta_j$
(counted with multiplicity, as in item 1).  Jensen's formula on $|t|\le r$
gives
\[
  \log\max_{|t|=r}|F^*|\ \ge\ \log|f_D|+\sum_{j:\rho_j>1/r}\log(r\rho_j)
  \ \ge\ \log|f_D|+\sum_{j}\log\bigl(r\max(1,\rho_j)\bigr),
\]
since each omitted or added term is $\le0$ on the right.  On the other side
$\max_{|t|=r}|F^*|\le|f_D|+b_1(F)\sum_{d\ge1}r^d
=|f_D|+b_1(F)\,r/(1-r)$.  Rearranging gives the first bound; for
$r=N/(N+1)$, $(1-r)/r=1/N$ and $r^N=(1+1/N)^{-N}\ge e^{-1}$.
\end{proof}

For real roots $r_j>1$ item 1 is the row $k=1$ of \cite[Theorem~3.2]{coefficient-mass}.  Both items
are of the size of the Mahler measure: they charge each root once, as
$\log\rho_j$, and do not give a second large coefficient.  At arbitrary
angles nothing more is true.

\begin{proposition}[No angle-uniform rows beyond the first]
\label{prop:noangle}
\leavevmode
\begin{enumerate}
\item Let $M\ge3$ be odd, $\rho>1$, and $0<\delta<\pi/2$.  The monic
polynomial $x^{2M}+\rho^{2M}$ is divisible by the product of the
$K=2\lfloor\delta M/\pi\rfloor+1$ conjugate pairs of modulus $\rho$ with
arguments $\pm(\frac\pi2+t\frac\pi M)$ for the integers $|t|\le\delta M/\pi$, all within
$\delta$ of the imaginary axis; yet
\[
  b_2(x^{2M}+\rho^{2M})=0,\qquad
  \Lambda(x^{2M}+\rho^{2M})=2M\log\rho\le\frac\pi\delta(K+1)\log\rho .
\]
\item For $m\ge2$ and real $a\ge1$, the polynomial $x^{2m}+x+a$ has no real
root, so its roots form $m$ conjugate pairs; all of them have moduli in
$[\rho/3,3\rho]$, where $\rho=a^{1/(2m)}$, and $b_2=1$,
$\Lambda=\log a=2m\log\rho$.
\end{enumerate}
In particular, no bound $b_2(F)\ge\phi(\rho_1,\ldots,\rho_K)>0$ in terms of
the moduli of the prescribed pairs alone, not even for pairs within a fixed
angle $\delta$ of the imaginary axis, and no bound
$\Lambda(F)\ge c\,K^2\log\rho$ for $K$ pairs of modulus $\rho$ within a
fixed angle $\delta$, can hold for all monic multiples.  Item~2 gives
the same linear mass without equal moduli or symmetric angles: $m$ pairs
with moduli within a factor $3$ of $\rho$, and a multiple with $b_2=1$ and
$\Lambda=2m\log\rho$.
\end{proposition}

\begin{proof}
1.  The roots of $x^{2M}+\rho^{2M}$ are $\rho e^{i\pi(2l+1)/(2M)}$; for
$M$ odd, $l=(M-1)/2+t$ gives the argument $\frac\pi2+t\frac\pi M$, and the
conjugate root has the negated argument.  The integers $t$ with
$|t|\le\delta M/\pi$ number $K$.  The only nonleading nonzero coefficient
is $\rho^{2M}$, so $b_2=0$ and $\Lambda=2M\log\rho$.  Finally
$K\ge2\delta M/\pi-1$ gives $2M\le\pi(K+1)/\delta$.

2.  For real $x$ with $|x|\ge1$, $x^{2m}+x\ge|x|^{2m}-|x|\ge0$, and for
$|x|<1$, $x^{2m}+x>-1$; so $x^{2m}+x+a>0$ when $a\ge1$.  The coefficients are
$1,1,a$.  The point $(1,\log1)=(1,0)$ lies on or below the chord from
$(0,\log a)$ to $(2m,0)$, whose height at $1$ is $(1-\frac1{2m})\log a\ge0$,
so the only tropical root (defined before Theorem~\ref{thm:archnewton}) is
$\log\rho=(\log a)/(2m)$, and Theorem~\ref{thm:archnewton}, item~2, puts every
root in $[\rho/3,3\rho]$.

For the final sentence, item 1 with $\delta M\ge\pi$, so $K\ge3$, has $b_2=0$ while every
$\rho_j=\rho$ is fixed, and letting $M\to\infty$ at fixed $\delta$ gives
$K\to\infty$ with $\Lambda=O_\delta(K\log\rho)$.
\end{proof}

What the imaginary axis provides in Corollary~\ref{cor:imag} is therefore
not closeness of the angles but exact membership, which makes the
prescribed factor a polynomial in $x^2$.  Separated moduli, on the other
hand, do help, through the Archimedean Newton polygon.  For a nonzero
$F=\sum_if_ix^i\in\mathbb C[x]$ of degree $D$ and order $e$ at $0$, let
$\varphi$ be the least concave function on $[e,D]$ with
$\varphi(i)\ge\log|f_i|$ whenever $f_i\ne0$.  Its vertices
$D=v_0>v_1>\cdots>v_J=e$ are among these points, so $\varphi(v_t)=\log|f_{v_t}|$.
The edge joining $v_s$ to $v_{s-1}$ has length $\ell_s=v_{s-1}-v_s\ge1$ and
slope $-\sigma_s$; by concavity $\sigma_1>\sigma_2>\cdots>\sigma_J$, and
\[
  \log|f_{v_t}|=\log|f_D|+\sum_{s\le t}\ell_s\sigma_s\qquad(0\le t\le J).
\]
We call $\sigma_1,\ldots,\sigma_J$ the \emph{tropical roots} of $F$.  For
$u\in\R$ the vertex $v_s$ with $\sigma_{s+1}<u<\sigma_s$ (where
$\sigma_0=+\infty$ and $\sigma_{J+1}=-\infty$) is the one maximizing
$\log|f_i|+iu$; we call it \emph{active} at $u$.

\begin{theorem}[Archimedean Newton polygon]\label{thm:archnewton}
Let $F\in\mathbb C[x]$ be nonzero, with tropical roots $\sigma_1>\cdots>\sigma_J$:
the edges of the least concave majorant $\varphi$ of the points
$(i,\log|f_i|)$, $f_i\ne0$, have slopes $-\sigma_1,\ldots,-\sigma_J$.  For
$u\in\R$ other than a tropical root, the vertex \emph{active} at $u$ is the
position maximizing $\log|f_i|+iu$.
\begin{enumerate}
\item Let $u\in\R$ satisfy $|u-\sigma_s|>\log3$ for every $s$, let $v$ be the
vertex active at $u$, and $r=e^u$.  Then
$|f_v|r^v>\sum_{i\ne v}|f_i|r^i$.  Hence $F$ has no root on $|z|=r$ and
exactly $v$ roots, counted with multiplicity, in $|z|<r$.
\item Every nonzero root $\beta$ of $F$ satisfies
$\bigl|\log|\beta|-\sigma_s\bigr|\le\log3$ for some $s$.
\end{enumerate}
\end{theorem}

\begin{proof}
1.  Let $v=v_s$, so $\sigma_{s+1}<u<\sigma_s$, and put $d_+=\sigma_s-u$ and
$d_-=u-\sigma_{s+1}$, both larger than $\log3$ (and infinite when $s=0$,
respectively $s=J$).  A concave function lies below the line through any of
its edges, so for $i>v$ the edge of slope $-\sigma_s$ gives
$\log|f_i|\le\varphi(i)\le\varphi(v)-\sigma_s(i-v)$, that is,
$|f_i|r^i\le|f_v|r^ve^{-(i-v)d_+}$; likewise the edge of slope $-\sigma_{s+1}$
gives $|f_i|r^i\le|f_v|r^ve^{-(v-i)d_-}$ for $i<v$.  Summing over the
distinct positive integers $|i-v|$,
\[
  \sum_{i\ne v}|f_i|r^i\ \le\ |f_v|r^v\Bigl(\frac1{e^{d_+}-1}+\frac1{e^{d_-}-1}\Bigr)
  \ <\ |f_v|r^v\Bigl(\frac12+\frac12\Bigr).
\]
So $|F(z)|\ge|f_v|r^v-\sum_{i\ne v}|f_i|r^i>0$ on $|z|=r$, and by Rouch\'e's
theorem $F$ and $f_vz^v$ have the same number of zeros in $|z|<r$, namely $v$.

2.  If $|\log|\beta|-\sigma_s|>\log3$ for every $s$, item~1 at
$u=\log|\beta|$ shows $F(\beta)\ne0$.
\end{proof}

Item~2 is the Archimedean form of the Newton polygon theorem: up to the
factor $3$, the moduli of the roots can only sit at the tropical roots.
Bounds of this kind are classical: they go back to Hadamard and to
Ostrowski's analysis of the Graeffe method~\cite{ostrowski}; item~1 is
Pellet's theorem~\cite{marden} at the active vertex; and the tropical roots
and log-majorization of root moduli by them are developed
in~\cite{akian-gaubert-sharify}.  We include the short proof to fix the
constant $3$.
Unlike the non-Archimedean case, the roots of modulus near $e^{\sigma_s}$ need not
number $\ell_s$; $(x-1)^D$ has all its roots on $|z|=1$ but
tropical roots $\log\frac{i+1}{D-i}$ spread over $[-\log D,\log D]$.
Separated roots still force separated vertices, which is all the mass needs.

\begin{corollary}[Separated moduli]\label{cor:separated}
Let $F\in\mathbb C[x]$ be nonzero of degree $D$, and let
$\rho_1\le\cdots\le\rho_K$ be moduli of nonzero roots of $F$.
\begin{enumerate}
\item If $\rho_1\ge3$ and $\rho_{j+1}>9\rho_j$ for $1\le j<K$, there are
positions $y_1<\cdots<y_K<D$ with
\[
  |f_{y_j}|\ \ge\ |f_D|\prod_{i=j}^K\frac{\rho_i}3 .
\]
Hence $b_k(F)\ge|f_D|\prod_{i=k}^K(\rho_i/3)$ for $1\le k\le K$, and, if
$|f_D|\ge1$,
\[
  \Lambda(F)\ \ge\ (K+1)\log|f_D|+\sum_{j=1}^Kj\log\frac{\rho_j}3 .
\]
\item If moreover each $\rho_j$ is the modulus of two roots of $F$, counted
with multiplicity (for instance $\beta_j$ and $\bar\beta_j$ with $F$ real and
$\beta_j\notin\R$), and $\rho_1>27$ and $\rho_{j+1}>729\rho_j$, then
$b_k(F)\ge|f_D|\prod_{i=k}^K(\rho_i^2/729)$ for $1\le k\le K$, and, if
$|f_D|\ge1$,
\[
  \Lambda(F)\ \ge\ (K+1)\log|f_D|+\sum_{j=1}^Kj\log\frac{\rho_j^2}{729} .
\]
\end{enumerate}
\end{corollary}

\begin{proof}
1.  By Theorem~\ref{thm:archnewton}, item~2, choose for each $j$ an index
$s(j)$ with $|\sigma_{s(j)}-\log\rho_j|\le\log3$.  Since
$\log\rho_{j+1}-\log\rho_j>2\log3$,
\[
  \sigma_{s(j+1)}\ge\log\rho_{j+1}-\log3>\log\rho_j+\log3\ge\sigma_{s(j)},
\]
so $s(1)>s(2)>\cdots>s(K)\ge1$, and $y_j=v_{s(j)}$ are increasing positions
below $D$.  In
$\log|f_{y_j}|=\log|f_D|+\sum_{s\le s(j)}\ell_s\sigma_s$ every term has
$\ell_s\ge1$ and $\sigma_s\ge\sigma_{s(j)}\ge\log(\rho_j/3)\ge0$; keeping only
the terms $s=s(i)$, $i\ge j$, gives
$\log|f_{y_j}|\ge\log|f_D|+\sum_{i\ge j}\log(\rho_i/3)$.  Every factor
$\rho_i/3$ is at least $1$, so the $k$ positions $y_1,\ldots,y_k$ all carry
magnitudes at least $|f_D|\prod_{i\ge k}(\rho_i/3)$, which bounds $b_k(F)$.
For the mass, $\Lambda(F)\ge\log^+|f_D|+\sum_j\log^+|f_{y_j}|$, and
$\sum_j\sum_{i\ge j}\log(\rho_i/3)=\sum_ii\log(\rho_i/3)$.

2.  Fix $c$ with $3\log3<c<\log\rho_1$ and $2c<\log(\rho_{j+1}/\rho_j)$ for
every $j$, and let $W_j=[\log\rho_j-c,\log\rho_j+c]$.  These windows are
disjoint, lie in $(0,\infty)$, and $W_{j+1}$ lies above $W_j$.  We claim that
the tropical roots in $W_j$ have total edge length
$\sum_{\sigma_s\in W_j}\ell_s\ge2$.  By Theorem~\ref{thm:archnewton}, item~2,
$W_j$ contains a tropical root $\sigma_s$ with $|\sigma_s-\log\rho_j|\le\log3$.
If it contains another one, the claim holds.  Otherwise take
$0<\varepsilon<c-3\log3$ and $u_\pm=\sigma_s\pm(\log3+\varepsilon)$.  Both are
farther than $\log3$ from $\sigma_s$, and every other tropical root lies
outside $W_j$, at distance more than $c-(2\log3+\varepsilon)>\log3$ from
$u_\pm$, because $|u_\pm-\log\rho_j|\le2\log3+\varepsilon$.  No tropical root
other than $\sigma_s$ lies in $[u_-,u_+]\subset W_j$, so the vertices active
at $u_+$ and $u_-$ are $v_{s-1}$ and $v_s$, and by
Theorem~\ref{thm:archnewton}, item~1, $F$ has exactly
$v_{s-1}-v_s=\ell_s$ roots in $e^{u_-}<|z|<e^{u_+}$.  The two roots of modulus
$\rho_j$ are among them, as $|\log\rho_j-\sigma_s|\le\log3$; so $\ell_s\ge2$.

Now let $\sigma_{s(j)}$ be the smallest tropical root in $W_j$ and
$y_j=v_{s(j)}$.  In $\log|f_{y_j}|=\log|f_D|+\sum_{s\le s(j)}\ell_s\sigma_s$
every $\sigma_s$ is at least $\sigma_{s(j)}>0$, and the sum contains every
tropical root in $W_i$ for $i\ge j$; each of these is at least
$\log\rho_i-c$.  By the claim,
$\log|f_{y_j}|\ge\log|f_D|+2\sum_{i\ge j}(\log\rho_i-c)$.  As in item~1,
$b_k(F)\ge|f_D|\prod_{i\ge k}(\rho_i^2e^{-2c})$ and
$\Lambda(F)\ge(K+1)\log|f_D|+\sum_jj\log(\rho_j^2e^{-2c})$.  These hold for
every $c$ in a nonempty interval with infimum $3\log3$; letting
$c\downarrow3\log3$ gives the stated bounds, as $e^{6\log3}=729$.

\end{proof}

So when the prescribed moduli grow by a factor $9$ at each step, every
complex multiple with $|f_D|\ge1$ pays $\sum_jj\log(\rho_j/3)$, one real root $\rho_j/3$ per
root, at every angle; with a factor $729$ a conjugate pair pays
$\log(\rho_j^2/729)$, which is the charge $\log(c_j^2-1)$ of
Corollary~\ref{cor:imag} up to the constant.  Proposition~\ref{prop:noangle}
shows that some separation is necessary; the constants $3$, $9$ and $729$ are
not optimized.

\section{Several annuli}\label{sec:annuli}

At complex roots in separated annuli the sharp worst-case mass bound charges
a root once for every annulus at or below it, and the cross terms between
annuli that \eqref{eq:order} has at real roots fail.  For real roots the
mass bound \eqref{eq:order} charges the $i$-th smallest
root $i$ times: once for itself and once for every smaller root.  Group $K_s$
real roots near $T_s$, with $T_1\ll\cdots\ll T_S$, and it reads
$\sum_s\bigl(K_1+\cdots+K_{s-1}+\frac{K_s+1}2\bigr)K_s\log T_s$ up to lower
order terms, with a square term $\frac12K_s^2\log T_s$ inside each group and a
cross term $K_uK_s\log T_s$ between the groups $u<s$.  At complex roots the
square term is false (Proposition~\ref{prop:noangle}), and so are the cross
terms: when the moduli of the $K_s$ roots
of the $s$-th group lie in $[T_s,U_s]$, with $T_1>3$ and $T_s>162U_{s-1}$, every multiple with $|f_D|\ge1$ has mass at least
$\sum_ssK_s\log(T_s/2)-S\log2$, and for suitable roots this is
attained to leading order (the ratio tends to $1$) as $\min_sT_s\to\infty$
(Theorem~\ref{thm:fixedgap}, Proposition~\ref{prop:annulisharp}).
Corollary~\ref{cor:separated} is essentially the case $K_s=1$.

For a nonzero $F$ and $r>0$ we call a position $y$ maximizing $|f_i|r^i$ a
\emph{central index} at $r$, as in Wiman--Valiron theory~\cite{hayman}.  In the notation of Theorem~\ref{thm:archnewton},
the central index at $r=e^u$ is the vertex active at $u$ when $u$ is not a
tropical root; at a tropical root it is any position $i$ of the corresponding
edge with $\log|f_i|=\varphi(i)$, its two vertices included.  Every central
index $y$ at $e^u$ satisfies
\begin{equation}\label{eq:centralabove}
  D-y\ \ge\ \sum_{\sigma_s>u}\ell_s .
\end{equation}
Indeed let $s$ be the largest index with
$\sigma_s>u$ ($s=0$ if there is none, and then $v_0=D$ and there is nothing
to prove); then $D-v_s=\sum_{t\le s}\ell_t$, and for $i>v_s$ the edge of
slope $-\sigma_s$ gives
$\log|f_i|+iu\le\log|f_{v_s}|+v_su-(\sigma_s-u)(i-v_s)<\log|f_{v_s}|+v_su$,
so $y\le v_s$.

\begin{lemma}[Central index]\label{lem:central}
Let $F\in\mathbb C[x]$ be nonzero of degree $D$, let $r>0$, and let $y$ be a
central index of $F$ at $r$.  For every sub-multiset $A$ of the zeros of $F$,
counted with multiplicity,
\[
  |f_y|\ >\ \frac{|f_D|}2\,r^{D-y-|A|}\prod_{\beta\in A}\max\Bigl(r,\frac{|\beta|}2\Bigr).
\]
In particular, if $r\ge1$ then $|f_y|>\frac{|f_D|}2\prod_{\beta\in A}\frac{|\beta|}{2r}$,
and if moreover $D-y\ge|A|$ then $|f_y|>\frac{|f_D|}2\prod_{\beta\in A}\frac{|\beta|}2$.
\end{lemma}

\begin{proof}
Let $\beta_1,\ldots,\beta_D$ be the zeros of $F$ with multiplicity.  Jensen's
formula on $|z|=2r$ and the choice of $y$ give
\[
  |f_D|\prod_{l=1}^D\max(2r,|\beta_l|)\ \le\ \max_{|z|=2r}|F(z)|
  \ \le\ \sum_{i=0}^D|f_i|r^i2^i\ \le\ |f_y|r^y\sum_{i=0}^D2^i\ <\ |f_y|r^y2^{D+1}.
\]
Each of the $D-|A|$ factors $\max(2r,|\beta_l|)$ with $\beta_l\notin A$ is at
least $2r$, so
$|f_y|>|f_D|2^{-D-1}r^{-y}(2r)^{D-|A|}\prod_{\beta\in A}\max(2r,|\beta|)$,
which is the claim.  For $r\ge1$, $r^{D-y-|A|}\prod_A\max(r,|\beta|/2)\ge
r^{D-y}\prod_A|\beta|/(2r)\ge\prod_A|\beta|/(2r)$; if $D-y\ge|A|$, then
$r^{D-y-|A|}\ge1$ and $\max(r,|\beta|/2)\ge|\beta|/2$.
\end{proof}

At $r=1$ the exponent $D-y-|A|$ plays no role: the central coefficient at
radius $1$ carries every zero, as $\max(1,|\beta|/2)$, like the first row of
Proposition~\ref{prop:rowone}.  At a larger radius the lemma charges
$|\beta|/2$ only for as many zeros as there are positions above the central
index.  The next lemma supplies such positions.

\begin{lemma}[Counting at a distance]\label{lem:tropcount}
Let $F\in\mathbb C[x]$ be nonzero of degree $D$, let $A$ be a sub-multiset of
the zeros of $F$, counted with multiplicity, with $n=|A|\ge1$ and $|\beta|\ge
T>0$ for $\beta\in A$, and let $0<r<T$ and $\theta=r/T$.  If $0\le\kappa\le n$ and
\[
  (1-\theta)^{-n}\sum_{m=n-\kappa+1}^n\binom nm\theta^m\ <\ 1,
\]
then every central index $y$ of $F$ at $r$ satisfies $D-y\ge\kappa$.  In
particular $D-y\ge n$ if $T\ge3nr$, and $D-y\ge n/2$ if $T\ge6r$.
\end{lemma}

\begin{proof}
Let $H=\prod_{\beta\in A}(x-\beta)=\sum_mh_mx^m$ and $P=F/H=\sum_jp_jx^j$, a
polynomial of degree $D-n$.  As $h_m=h_0(-1)^me_m(1/\beta)$ and $h_0\ne0$,
$|h_m|r^m\le|h_0|\binom nm\theta^m$.  In $\mathbb C[[x]]$,
$1/H=h_0^{-1}\prod_\beta(1-x/\beta)^{-1}=\sum_mq_mx^m$ with
$|q_m|\le|h_0|^{-1}\binom{n+m-1}mT^{-m}$, and $P=F\cdot(1/H)$, so
$p_j=\sum_{m\le j}f_{j-m}q_m$.  With $M=\max_{i\le D-n}|f_i|r^i$, for
$j\le D-n$,
\[
  |p_j|r^j\ \le\ \sum_m|f_{j-m}|r^{j-m}\,|q_m|r^m
  \ \le\ \frac M{|h_0|}\sum_m\binom{n+m-1}m\theta^m\ =\ \frac M{|h_0|}(1-\theta)^{-n};
\]
in particular $M>0$, as $P\ne0$.  For $i>D-\kappa$, $f_i=\sum_jp_jh_{i-j}$
over $j\le D-n$, so $i-j>n-\kappa$ and
\[
  |f_i|r^i\ \le\ \sum_j|p_j|r^j\,|h_{i-j}|r^{i-j}
  \ \le\ M(1-\theta)^{-n}\sum_{m>n-\kappa}\binom nm\theta^m\ <\ M .
\]
So no $i>D-\kappa$ is a central index.  If $T\ge3nr$ and $\kappa=n$, the sum
is $(1+\theta)^n-1\le e^{1/3}-1<\frac23\le1-n\theta\le(1-\theta)^n$.  If $T\ge6r$ and
$\kappa=\lceil n/2\rceil$, every $m>n-\kappa$ has $m\ge(n+1)/2$, so the left
side is at most $(1-\theta)^{-n}2^n\theta^{(n+1)/2}\le6^{-1/2}(12/(5\sqrt6))^n<1$.
\end{proof}

The factor $3n$ is sharp up to the $3$: $(x-t)^K$ has $K$ zeros of modulus
$t$, and at $r=t/c$ its terms are $t^K\binom Kic^{-i}$, so for $c<K$ the term
$i=1$ beats the term $i=0$ and the central index is positive.  At a constant
distance only a fraction of the count survives: for $c\ge1$ the central
index of $(x-t)^K$ at $t/c$ is about $K/(c+1)$, with only about $cK/(c+1)$
positions above it.

\begin{theorem}[Several annuli]\label{thm:annuli}
Let $F\in\mathbb C[x]$ be nonzero of degree $D$, let $S\ge1$, and let
$A_1,\ldots,A_S$ be nonempty sub-multisets of the zeros of $F$, counted with
multiplicity, the moduli in $A_s$ lying in $[T_s,U_s]$, where
\[
  T_1>3,\qquad T_a>9U_{a-1}\quad(2\le a\le S).
\]
Put $K_s=|A_s|$, $n_a=K_a+\cdots+K_S$ and $A_{\ge a}=A_a\cup\cdots\cup A_S$.
Let $r_1=1$ and $k_1=n_1$, and for $2\le a\le S$ choose $r_a$ with
$3U_{a-1}<r_a<T_a/3$, let $\kappa_a$ be the largest integer with
$3\kappa_ar_a\le T_a$, and put $k_a=\min\bigl(n_a,\max(\kappa_a,S-a+1)\bigr)$.
Then there are positions $y_1<y_2<\cdots<y_S<D$ with
\[
  |f_{y_a}|\ >\ B_a=\frac{|f_D|}2\,r_a^{k_a-n_a}\prod_{\beta\in A_{\ge a}}\frac{|\beta|}2
  \qquad(1\le a\le S).
\]
Hence $b_k(F)>\min_{a\le k}B_a$ for $1\le k\le S$, and, if $|f_D|\ge1$,
\[
  \Lambda(F)\ \ge\ (S+1)\log|f_D|+\sum_{s=1}^Ss\sum_{\beta\in A_s}\log\frac{|\beta|}2
  -\sum_{a=2}^S(n_a-k_a)\log r_a-S\log2 .
\]
\end{theorem}

\begin{proof}
Let $y_a$ be a central index of $F$ at $r_a$.  Since $r_a<T_a/3$,
$\kappa_a\ge1$.

\emph{Order.}  For $1\le j\le S$ pick $\beta\in A_j$.  By
Theorem~\ref{thm:archnewton}, item~2, some tropical root $\sigma^{(j)}$
satisfies $|\sigma^{(j)}-\log|\beta||\le\log3$, so
$\log(T_j/3)\le\sigma^{(j)}\le\log(3U_j)$.  These intervals are disjoint, as
$T_{j+1}>9U_j$, so the $\sigma^{(j)}$ are distinct.  For every $a$,
$\log r_a<\log(T_a/3)\le\sigma^{(a)}$ (for $a=1$ because $T_1>3$), and for
$a<S$ also $\sigma^{(a)}\le\log(3U_a)<\log r_{a+1}$.  Let
$\sigma^{(a)}=\sigma_q$ in the numbering of Theorem~\ref{thm:archnewton}.  By the
proof of \eqref{eq:centralabove}, a central index at a
radius below $e^{\sigma_q}$ is at most the lower vertex $v_q$ of the edge of
slope $-\sigma_q$; symmetrically, for $i<v_{q-1}$ the same edge gives
$\log|f_i|+ip<\log|f_{v_{q-1}}|+v_{q-1}p$ when $p>\sigma_q$, so a central
index at a radius above $e^{\sigma_q}$ is at least the upper vertex
$v_{q-1}>v_q$.  Hence $y_a\le v_q<v_{q-1}\le y_{a+1}$ for $a<S$, and
$y_S\le v_q<v_0=D$ for $a=S$.

\emph{Counts.}  Let $a\ge2$.  The $S-a+1$ distinct tropical roots
$\sigma^{(a)},\ldots,\sigma^{(S)}$ exceed $\log r_a$, so
$D-y_a\ge S-a+1$ by \eqref{eq:centralabove}.  Let $n=\min(n_a,\kappa_a)$;
Lemma~\ref{lem:tropcount} with $A$ any $n$ of the zeros in $A_{\ge a}$, of
modulus at least $T_a\ge3nr_a$, gives $D-y_a\ge n$.  So $D-y_a\ge k_a$ for every $a\ge2$.

\emph{Values.}  Every $r_a\ge1$ (for $a\ge2$, $r_a>3U_{a-1}\ge3T_1>9$).
Lemma~\ref{lem:central} with $A=A_{\ge a}$ gives
$|f_{y_a}|>\frac{|f_D|}2r_a^{D-y_a-n_a}\prod_{A_{\ge a}}\max(r_a,|\beta|/2)$.
For $a=1$ the power of $r_1=1$ is $1$; for $a\ge2$ it is at least
$r_a^{k_a-n_a}$, since $D-y_a\ge k_a$; and $\max(r_a,|\beta|/2)\ge|\beta|/2$.
This is $|f_{y_a}|>B_a$.

The positions $y_1,\ldots,y_k$ are $k$ distinct nonleading positions with
magnitudes above $\min_{a\le k}B_a$, which bounds $b_k(F)$.  For the mass,
$\Lambda(F)\ge\log^+|f_D|+\sum_a\log^+|f_{y_a}|\ge\log|f_D|+\sum_a\log B_a$,
and $\sum_a\sum_{\beta\in A_{\ge a}}=\sum_ss\sum_{\beta\in A_s}$; the term
$a=1$ has no power of $r_1$.
\end{proof}

\begin{corollary}[Annuli]\label{cor:annuli}
In the setting of Theorem~\ref{thm:annuli}:
\begin{enumerate}
\item If moreover $T_a>9n_aU_{a-1}$ for $2\le a\le S$, then
$b_k(F)>\frac{|f_D|}2\prod_{\beta\in A_{\ge k}}\frac{|\beta|}2$ for $1\le k\le S$,
and, if $|f_D|\ge1$,
\[
  \Lambda(F)\ \ge\ (S+1)\log|f_D|+\sum_{s=1}^Ss\sum_{\beta\in A_s}\log\frac{|\beta|}2-S\log2 .
\]
\item Without the extra hypothesis, if $|f_D|\ge1$,
\[
  \Lambda(F)\ \ge\ (S+1)\log|f_D|+\sum_{\beta\in A_{\ge1}}\log\frac{|\beta|}2
  +\sum_{a=2}^S\sum_{\beta\in A_{\ge a}}\log\frac{|\beta|}{6U_{a-1}}-S\log2 .
\]
\end{enumerate}
\end{corollary}

\begin{proof}
1.  Choose $r_a$ with $3U_{a-1}<r_a<T_a/(3n_a)$.  Then $r_a<T_a/3$ and
$\kappa_a\ge n_a$, so $k_a=n_a$ and
$B_a=\frac{|f_D|}2\prod_{A_{\ge a}}\frac{|\beta|}2$.  Every factor
$|\beta|/2$ exceeds $1$, so $B_a$ decreases with $a$ and
$\min_{a\le k}B_a=B_k$.

2.  As $k_a\ge0$, $(n_a-k_a)\log r_a\le n_a\log r_a$, and
$\sum_{A_{\ge a}}\log\frac{|\beta|}2-n_a\log r_a=\sum_{A_{\ge a}}\log\frac{|\beta|}{2r_a}$.
The bound holds for every admissible $r_a$; let $r_a\downarrow3U_{a-1}$.
\end{proof}

Item~1 charges a root of $A_s$ once for each annulus at or below it, $s$
times in all, at the price of a separation that grows linearly with the
number of roots above; by the remark after Lemma~\ref{lem:tropcount}, a
central index at $r_a$ can have fewer than $n_a$ positions above it when
$T_a/r_a<n_a$.  Item~2 needs only the factor $9$, but past the first
annulus it charges ratios $|\beta|/U_{a-1}$ instead of moduli.  With one root
per annulus, $k_a=n_a=S-a+1$ at the separation $9$ of the theorem, which then
gives Corollary~\ref{cor:separated}, item~1, when $\rho_1>3$ (the corollary
allows $\rho_1=3$, the theorem needs $T_1>3$), with $\rho_j/2$ in place of
$\rho_j/3$ and an extra factor $\frac12$ in each row.  For the mass alone a
fixed separation suffices: where one central index falls short, two of them
share the charge.

\begin{theorem}[Annuli at a fixed separation]\label{thm:fixedgap}
Let $F\in\mathbb C[x]$ be nonzero of degree $D$ with $|f_D|\ge1$, let
$S\ge1$, and let $A_1,\ldots,A_S$ be nonempty sub-multisets of the zeros of
$F$, counted with multiplicity, the moduli in $A_s$ lying in $[T_s,U_s]$,
where
\[
  T_1>3,\qquad T_a>162\,U_{a-1}\quad(2\le a\le S).
\]
Then
\[
  \Lambda(F)\ \ge\ (S+1)\log|f_D|+\sum_{s=1}^Ss\sum_{\beta\in A_s}\log\frac{|\beta|}2-S\log2 .
\]
\end{theorem}

\begin{proof}
Use the notation of Theorem~\ref{thm:annuli} and put
$Z_a=\sum_{\beta\in A_{\ge a}}\log(|\beta|/2)$, so that
$\sum_aZ_a=\sum_ss\sum_{\beta\in A_s}\log(|\beta|/2)$ and
$Z_a\ge n_a\log(T_a/2)$.  The Order step of that proof uses only
$T_1>3$ and $T_{a}>9U_{a-1}$; it gives tropical roots
$\sigma^{(1)}<\cdots<\sigma^{(S)}$ with
$\log(T_a/3)\le\sigma^{(a)}\le\log(3U_a)$, and shows that a central index at
a radius below $e^{\sigma^{(a)}}$ lies below every central index at a radius
above it.  For $2\le a\le S$ choose
$u_a^-<u_a^+$, neither a tropical root, with
\[
  \log(3U_{a-1})<u_a^-,\qquad u_a^-+2\log3<u_a^+\le\log(T_a/6);
\]
this is possible as $T_a/6>9\cdot3U_{a-1}$.  Let $Y_a$ be the set of the
vertices active at $u_a^-$ and at $u_a^+$, and let $y_1$ be a central index
at $1$.  As $\sigma^{(a-1)}<u_a^-<u_a^+<\sigma^{(a)}$ and $0<\sigma^{(1)}$,
$y_1$ lies below every position in $Y_2$, every
position in $Y_a$ below every position in $Y_{a+1}$, and every position in
$Y_S$ below $D$ ($y_1$ below $D$ if $S=1$).  So the positions in $\{y_1\}\cup Y_2\cup\cdots\cup Y_S$
are distinct and nonleading.

By Lemma~\ref{lem:central} at $r=1$ with $A=A_{\ge1}$,
$\log|f_{y_1}|>\log(|f_D|/2)+Z_1$.  Let $a\ge2$ and let $y\in Y_a$ be active
at $u\in\{u_a^-,u_a^+\}$, so $9<e^u\le T_a/6$.  Every $\beta\in A_{\ge a}$ has
$|\beta|/2>e^u$, so Lemma~\ref{lem:central} with $A=A_{\ge a}$ gives
$\log|f_y|>\log(|f_D|/2)+Z_a-(n_a-D+y)u$, and
Lemma~\ref{lem:tropcount} with $A=A_{\ge a}$, $T=T_a\ge6e^u$, gives
$D-y\ge n_a/2$.

If $Y_a=\{y\}$, no tropical root lies in $[u_a^-,u_a^+]$ (one would separate
the two vertices, as in the Order step), so the midpoint $u$
of this interval is farther than $\log3$ from every tropical root, and $y$ is
the vertex active at $u$.  By Theorem~\ref{thm:archnewton}, item~1, $F$ has
exactly $D-y$ zeros in $|z|>e^u$; the $n_a$ zeros of $A_{\ge a}$ are among
them, so $D-y\ge n_a$ and $\log|f_y|>\log(|f_D|/2)+Z_a$.  Otherwise $Y_a$
has two positions, each with $n_a-D+y\le n_a/2$ at a level
$u\le u_a^+\le\log(T_a/6)$, and their logarithms add up to more than
\[
  2\log\frac{|f_D|}2+2Z_a-n_a\log\frac{T_a}6
  \ \ge\ 2\log\frac{|f_D|}2+Z_a+n_a\log3
  \ \ge\ \log\frac{|f_D|}2+Z_a ,
\]
as $\log|f_D|\ge0$ and $n_a\log3>\log2$.  In both cases
$\sum_{y\in Y_a}\log^+|f_y|\ge\log|f_D|+Z_a-\log2$, and
\[
  \Lambda(F)\ \ge\ \log^+|f_D|+\log^+|f_{y_1}|+\sum_{a=2}^S\sum_{y\in Y_a}\log^+|f_y|
  \ \ge\ (S+1)\log|f_D|+\sum_{a=1}^SZ_a-S\log2 .
\]
\end{proof}

Theorem~\ref{thm:fixedgap} bounds the mass only: below the $a$-th annulus
the charge of the roots in the $a$-th and higher annuli may fall on two
coefficients rather than one, and whether the rows of
Corollary~\ref{cor:annuli}, item~1, hold at a fixed separation is open.  The
factor $162$ is not optimized.  The mass bound is sharp to leading order, the
coefficient $s$ of each $K_s\log T_s$ being best possible, as
$\min_sT_s\to\infty$ (Proposition~\ref{prop:annulisharp} has
$\Lambda=\sum_ssN_s\log t_s$ against the bound
$\sum_ssN_s\log(t_s/6)-S\log2$), and the cross terms $K_uK_s\log T_s$ of the
real bound fail.

\begin{proposition}[Sharpness for annuli]\label{prop:annulisharp}
\leavevmode
\begin{enumerate}
\item Let $N_1,\ldots,N_S\ge1$ and let $2\le t_1<\cdots<t_S$ be integers with
$t_{a+1}>9t_a$.  Put $p_a=N_1+\cdots+N_a$ ($p_0=0$), $D=p_S$ and
\[
  F=\sum_{a=0}^Sc_ax^{p_a},\qquad c_a=\prod_{s>a}t_s^{N_s}.
\]
Then $F\in\mathbb Z[x]$ is monic, $F$ has exactly $N_s$ zeros, counted with
multiplicity, in $t_s/3\le|z|\le3t_s$ for each $s$ and no other zeros, and
\[
  \Lambda(F)=\sum_{s=1}^SsN_s\log t_s .
\]
If moreover $t_1>9$ and $t_a>1458\,t_{a-1}$, then
Theorem~\ref{thm:fixedgap} applies to $F$ with $A_s$ the zeros of modulus in
$[t_s/3,3t_s]$, and gives $\Lambda(F)\ge\sum_ssN_s\log(t_s/6)-S\log2$.
\item Let $0<\delta<\pi/2$, let $M_1\ne M_2$ be odd integers at least $3$,
and let $1<t_1<t_2$.  The monic real polynomial
$(x^{2M_1}+t_1^{2M_1})(x^{2M_2}+t_2^{2M_2})$ is divisible, for $s=1,2$, by
$K_s=2\lfloor\delta M_s/\pi\rfloor+1$ conjugate pairs of modulus $t_s$
within $\delta$ of the imaginary axis, and
\[
  \Lambda=4M_1\log t_1+4M_2\log t_2\ \le\
  \frac{2\pi}\delta\bigl((K_1+1)\log t_1+(K_2+1)\log t_2\bigr).
\]
\end{enumerate}
In particular, for no $c>0$ does $\Lambda(F)\ge c\,K_1K_2\log T_2$ hold for
all monic multiples $F$ of $K_1$ roots with moduli in $[T_1,9T_1]$ and $K_2$
roots with moduli in $[T_2,9T_2]$, however large $T_2/T_1$ is, nor with $K_s$
counting conjugate pairs within a fixed angle of the imaginary axis, as in
item~2.
\end{proposition}

\begin{proof}
1.  $F$ has the $S+1$ nonzero coefficients $c_a\in\mathbb Z$ at the positions
$p_a$, and $c_S=1$.  The slope from $(p_{a-1},\log c_{a-1})$ to
$(p_a,\log c_a)$ is $-\log t_a$, which decreases strictly with $a$, so these
points are the vertices of the least concave majorant, and the tropical roots
are the $\log t_a$, the edge of $\log t_a$ joining $p_{a-1}$ to $p_a$.  By
Theorem~\ref{thm:archnewton}, item~2, every zero lies in one of the closed
annuli $t_s/3\le|z|\le3t_s$, which are disjoint as $t_{s+1}>9t_s$.  For $0\le a\le S$
choose $u_a$ with $3t_a<e^{u_a}<t_{a+1}/3$ (with $t_0=0$ and $t_{S+1}=\infty$).
It is farther than $\log3$ from every tropical root, and the vertex active at
$u_a$ is $p_a$; by Theorem~\ref{thm:archnewton}, item~1, $F$ has exactly
$p_a$ zeros in $|z|<e^{u_a}$, that is, in the first $a$ annuli.  So the
$a$-th annulus holds $p_a-p_{a-1}=N_a$ zeros.  Every $c_a\ge1$, so
$\Lambda(F)=\sum_{a<S}\log c_a=\sum_{a<S}\sum_{s>a}N_s\log t_s
=\sum_ssN_s\log t_s$.

For the last claim take $T_s=t_s/3$ and $U_s=3t_s$: then $T_1>3$,
$T_a>162U_{a-1}$, $|A_s|=N_s$, $f_D=1$, and $|\beta|/2\ge t_s/6$ on $A_s$.

2.  By the proof of Proposition~\ref{prop:noangle}, item~1, each factor
$x^{2M_s}+t_s^{2M_s}$ has $K_s$ conjugate pairs of modulus $t_s$ within
$\delta$ of the imaginary axis, and $2M_s\le\pi(K_s+1)/\delta$.  The product
is $x^{2M_1+2M_2}+t_2^{2M_2}x^{2M_1}+t_1^{2M_1}x^{2M_2}+t_1^{2M_1}t_2^{2M_2}$,
four distinct positions as $M_1\ne M_2$, which gives $\Lambda$.

For the final sentence, item~1 with $S=2$, $N_1=N_2=N$, $t_1\ge9$ and $T_s=t_s/3$ (so $T_s\ge3$)
has $K_s=N$ zeros with moduli in $[T_s,9T_s]$ and
$\Lambda(F)=N\log t_1+2N\log t_2\le3N\log(3T_2)\le6N\log T_2$, which is below
$cN^2\log T_2$ once $N>6/c$.  Item~2 with $T_s=t_s$ has
$\Lambda\le\frac{2\pi}\delta(K_1+K_2+2)\log T_2$, below $cK_1K_2\log T_2$
once $M_1$ and $M_2$, hence $K_1$ and $K_2$, are large.
\end{proof}

So at complex roots in separated annuli the sharp worst-case analogue of
$\sum_ii\log r_i$ is $\sum_ssK_s\log T_s$: the cross term between annuli $u<s$ is $K_s\log T_s$,
linear in the number of roots above and independent of the number below.  For
the mass it is proved at the fixed separation $162$
(Theorem~\ref{thm:fixedgap}), and with ratios in place of moduli at the
separation $9$ (Corollary~\ref{cor:annuli}, item~2); the rows of item~1 are
proved at the separation $9n_a$.  Open: whether the rows hold
at a fixed separation, and the least fixed separation for the mass, in
particular whether $9$ suffices.

\section{Real roots of both signs}\label{sec:bothsigns}

\cite[Theorem~3.2]{coefficient-mass} needs its roots on one side of $0$.  Applied to $F(x)$
and to $F(-x)$ it charges a multiple with $p$ positive and $n$ negative
prescribed roots only for the larger of the two classes.  The mixed problem is
weaker: by the multisection transfer (the remark after
Corollary~\ref{cor:imag}, with $R=\prod_j(y-s_j^2)$), the $L=2q$ roots $\pm s_j$
of $\prod_j(x^2-s_j^2)$ are worth only the $q$ positive roots $s_j^2$, so at a
common modulus $s$ they cost about $\frac14L^2\log s$ instead of the
$\frac12L^2\log s$ of $L$ roots of one sign.  The operator $T_z$ of the
proof of \cite[Proposition~4.2]{coefficient-mass}, where Rolle's theorem is
used on $(0,\tfrac12]$ only, here used on the two half-lines separately,
gives the matching lower bound: every excluded coefficient costs one root of
each sign.

\begin{lemma}[Rolle on both half-lines]\label{lem:rolleside}
Let $h\in\R[t]$ be nonzero, let $z\ge1$ be an integer, and put
$T_zh=h-(t/z)h'$, so that $[t^d]T_zh=(1-d/z)[t^d]h$.  If the zeros of $h$ in
$(0,\infty)$, listed with multiplicity, are $\xi_1\le\cdots\le\xi_m$, $m\ge1$, then $T_zh$ has
at least $m-1$ zeros in $(0,\infty)$, counted with multiplicity, and $m-1$ of
them, $\eta_1\le\cdots\le\eta_{m-1}$, satisfy $\xi_j\le\eta_j\le\xi_{j+1}$.
The same holds in $(-\infty,0)$ for the
absolute values: zeros $|\xi_1|\le\cdots\le|\xi_m|$ of $h$ there give
$m-1$ zeros $|\eta_1|\le\cdots\le|\eta_{m-1}|$ of $T_zh$ there with
$|\xi_j|\le|\eta_j|\le|\xi_{j+1}|$.
\end{lemma}

\begin{proof}
For $t\ne0$, $T_zh=-(t^{z+1}/z)(t^{-z}h)'$.  On $I=(0,\infty)$ the function
$\varphi=t^{-z}h$ is real-analytic, not identically zero, and has the zeros of
$h$ with the same multiplicities.  Let $x_1<\cdots<x_q$ be the distinct zeros,
of multiplicities $\mu_i$.  Writing $\varphi=(t-x_i)^{\mu_i}\psi$ with
$\psi(x_i)\ne0$ shows that $\varphi'$ vanishes to order $\mu_i-1$ at $x_i$,
and Rolle's theorem gives a zero $y_i\in(x_i,x_{i+1})$ for $i<q$.  List
$\mu_1-1$ copies of $x_1$, then $y_1$, then $\mu_2-1$ copies of $x_2$, then
$y_2$, and so on: these are $m-1$ zeros of $\varphi'$ counted with
multiplicity, the $j$-th lies in $[\xi_j,\xi_{j+1}]$, and since
$-t^{z+1}/z$ has no zero in $I$ they are zeros of $T_zh$.  For
$(-\infty,0)$ apply this to $h(-t)$, using $(T_zh)(-t)=T_z\bigl(h(-t)\bigr)$.
\end{proof}

\begin{theorem}[Rows at real roots of both signs]\label{thm:bothsigns}
Let $F\in\mathbb C[x]$ be nonzero of degree $D$, and let
$a_1\ge\cdots\ge a_p>0$ and $c_1\ge\cdots\ge c_n>0$, repetitions allowed, with
$p+n\ge1$, be such that $\prod_{i}(x-a_i)\prod_i(x+c_i)$ divides $F$.  Let
$0<\rho<1$ with $\rho a_i>1$ and $\rho c_i>1$ for every $i$.  For
$1\le k\le\max(p,n)$ put
\[
  \Pi_k(\rho)=\prod_{i=k}^{p}(\rho a_i)\prod_{i=k}^{n}(\rho c_i),
\]
an empty product being $1$.  Then
\begin{equation}\label{eq:bothrowsbound}
  b_k(F)\ \ge\ |f_D|\,\frac{\Pi_k(\rho)-1}{(1-\rho)^{-k}-1}.
\end{equation}
\end{theorem}

\begin{proof}
First let $F$ be real.  Since $k\le\max(p,n)\le p+n\le D$, there are at least
$k-1$ positions below $D$.  Let $z_1<\cdots<z_{k-1}$ be the backward distances
$D-j$ of $k-1$ positions $j<D$ carrying the $k-1$ largest nonleading
magnitudes; they are distinct positive integers, so $z_i\ge i$.  Let
$F^*(t)=t^DF(1/t)=\sum_{d=0}^Df_{D-d}t^d$ and
\[
  g=T_{z_{k-1}}\cdots T_{z_1}F^*,\qquad
  [t^d]g=f_{D-d}\prod_{i=1}^{k-1}\Bigl(1-\frac d{z_i}\Bigr).
\]
So $g(0)=f_D\ne0$ and $[t^{z_i}]g=0$.  As in the proof of
Proposition~\ref{prop:rowone}, $(x-\alpha)\mid F$ gives $(1-\alpha t)\mid F^*$,
so $F^*$ has the zeros $1/a_1\le\cdots\le1/a_p$ in $(0,\infty)$ and the zeros
$-1/c_i$, of absolute values $1/c_1\le\cdots\le1/c_n$, in $(-\infty,0)$, with
multiplicity.  Every intermediate polynomial $T_{z_s}\cdots T_{z_1}F^*$ has
constant term $f_D\ne0$, so Lemma~\ref{lem:rolleside} applies at each step;
after $s$ steps the $j$-th smallest positive zero is at most $\eta_j$, hence
at most the $(j+1)$-th smallest positive zero of the previous stage, hence at most $1/a_{j+s}$.  So $g$ has positive
zeros $\eta_j\le1/a_{j+k-1}$, $1\le j\le p-k+1$, and negative zeros
$\eta'_j$ with $|\eta'_j|\le1/c_{j+k-1}$, $1\le j\le n-k+1$, all of modulus
less than $\rho$.  Jensen's formula for $g$ on $|t|\le\rho$ gives
\[
  \log|f_D|+\log\Pi_k(\rho)\ \le\
  \log|f_D|+\sum_{g(\zeta)=0,\ |\zeta|<\rho}\log\frac\rho{|\zeta|}\ \le\
  \log\max_{|t|=\rho}|g(t)|,
\]
because $\log(\rho/|\eta_j|)\ge\log(\rho a_{j+k-1})$, likewise for the
$\eta'_j$, and the remaining terms are positive.  On the other side, for
$d\ge1$ either $[t^d]g=0$ or the position $D-d$ is not among the $k-1$ largest,
so $|f_{D-d}|\le b_k(F)$, while
$\prod_i|1-d/z_i|\le\prod_{i=1}^{k-1}(1+d/i)=\binom{d+k-1}{k-1}$.  Hence
\[
  \max_{|t|=\rho}|g(t)|\ \le\ |f_D|+b_k(F)\sum_{d\ge1}\binom{d+k-1}{k-1}\rho^d
  =|f_D|+b_k(F)\bigl((1-\rho)^{-k}-1\bigr),
\]
and the two displays give \eqref{eq:bothrowsbound}.  For complex $F$ apply the
real case to $\operatorname{Re}(F/f_D)$, as in \cite[Corollary~3.3]{coefficient-mass}: it
is real of degree $D$ with leading coefficient $1$, divisible by the real
polynomial $\prod_i(x-a_i)\prod_i(x+c_i)$, and its coefficients are dominated
by those of $F/f_D$.
\end{proof}

The exclusions cost one root of each sign, and in the worst case the largest
one: $\Pi_k$ omits $a_1,\ldots,a_{k-1}$ and $c_1,\ldots,c_{k-1}$.  This is the
opposite of \cite[Theorem~3.2]{coefficient-mass}, which omits the smallest roots, so
\eqref{eq:bothrowsbound} is the right tool when the moduli are comparable,
and the mass bound it yields has the sharp constant $\frac14$
(Proposition~\ref{prop:bothsignssharp}).

\begin{corollary}[Mass at real roots of both signs]\label{cor:bothsignsmass}
Let $F\in\mathbb C[x]$ be nonzero with $|f_D|\ge1$, divisible by
$\prod_{i=1}^L(x-r_i)$ with $r_i$ real, repetitions allowed, and $|r_i|\ge R\ge4$.
Then for $1\le k\le\lceil L/2\rceil$,
\[
  b_k(F)\ \ge\ |f_D|\,\frac{(R/2)^{L-2k+2}-1}{2^k-1},
\]
and
\begin{equation}\label{eq:bothmass}
  \Lambda(F)\ \ge\ \Bigl(\Bigl\lceil\frac L2\Bigr\rceil+1\Bigr)\log|f_D|
  +\Bigl\lfloor\frac{(L+1)^2}4\Bigr\rfloor\log\frac R2
  -\frac12\Bigl\lceil\frac L2\Bigr\rceil\Bigl(\Bigl\lceil\frac L2\Bigr\rceil+3\Bigr)\log2 .
\end{equation}
\end{corollary}

\begin{proof}
Let $p$ and $n$ count the positive and negative $r_i$; then
$\lceil L/2\rceil\le\max(p,n)$.  Take $\rho=\frac12$ in
Theorem~\ref{thm:bothsigns}: every $\rho|r_i|\ge R/2\ge2$, and
$\Pi_k(\tfrac12)\ge(R/2)^{(p-k+1)^++(n-k+1)^+}\ge(R/2)^{L-2k+2}$.  This is the row
bound.  For $k\le\lceil L/2\rceil$ the exponent $L-2k+2$ is at least $1$, so
$X=(R/2)^{L-2k+2}\ge2$, $X-1\ge X/2$, and
$\log b_k(F)\ge\log|f_D|+(L-2k+2)\log(R/2)-(k+1)\log2$.  The leading coefficient
and the $\lceil L/2\rceil$ largest nonleading magnitudes, all nonzero, sit at
distinct positions, so $\Lambda(F)\ge\log^+|f_D|+\sum_k\log^+b_k(F)$ is at
least the sum of these lower bounds; finally
$\sum_{k=1}^{\lceil L/2\rceil}(L-2k+2)=\lfloor(L+1)^2/4\rfloor$ and
$\sum_{k=1}^{\lceil L/2\rceil}(k+1)=\frac12\lceil L/2\rceil(\lceil L/2\rceil+3)$.
\end{proof}

\begin{proposition}[Sharpness of the constant $\frac14$]\label{prop:bothsignssharp}
Let $m\ge1$, $R\ge4$, $\lambda\ge1$ and $s_1,\ldots,s_m\in[R,\lambda R]$.  The
polynomial $P=\prod_{j=1}^m(x^2-s_j^2)$ is a monic multiple of the $L=2m$ real
roots $\pm s_j$, all of modulus at least $R$, and
\[
  \Lambda(P)\ \le\ \Bigl\lfloor\frac{(L+1)^2}4\Bigr\rfloor\log(\lambda R)+m^2\log2 .
\]
So the coefficient $\lfloor(L+1)^2/4\rfloor\sim\frac14L^2$ of $\log R$ in
\eqref{eq:bothmass} cannot be increased as $R\to\infty$ with $m$ and
$\lambda$ fixed, where $m^2\log2$ is of lower order, while for roots of one sign
\cite[Corollary~3.4]{coefficient-mass} gives $\frac{L(L+1)}2\log(R-1)$.
\end{proposition}

\begin{proof}
The nonleading coefficients of $P$ are $\pm e_j(s_1^2,\ldots,s_m^2)$,
$1\le j\le m$, each between $1$ and $\binom mj(\lambda R)^{2j}$.  So
$\Lambda(P)\le\sum_j\log\binom mj+m(m+1)\log(\lambda R)$, and
$m(m+1)=\lfloor(2m+1)^2/4\rfloor$, $\sum_j\log\binom mj\le m^2\log2$.
\end{proof}

\section{Gaussian-integer roots}\label{sec:gaussian}

For roots in $\mathbb Z[i]$ and integer multiples there is one more piece of
information: an integer polynomial takes Gaussian-integer values at them,
and a nonzero Gaussian integer has modulus at least one.  This charges the
degree of a multiple against its mass (Theorem~\ref{thm:gausscharge}): an
integer multiple splits into blocks, and a block of large degree has large
mass.

Throughout, $\alpha_1,\ldots,\alpha_K\in\mathbb Z[i]$ are distinct, with
$\alpha_j=a_j+ic_j$ and $c_j\ge1$; put $\rho_j=|\alpha_j|$,
$\rho_{\min}=\min_j\rho_j$ and
\[
  P=\prod_{j=1}^K(x-\alpha_j)(x-\bar\alpha_j)
   =\prod_{j=1}^K\bigl(x^2-2a_jx+a_j^2+c_j^2\bigr)\in\mathbb Z[x],
\]
a monic polynomial with $2K$ distinct roots and $P(0)\ne0$.  For
$F=\sum_{i=0}^Df_ix^i$ and $0\le m\le D+1$ write
\[
  F_{<m}=\sum_{i<m}f_ix^i,\qquad G_m=\sum_{i\ge m}f_ix^{i-m},\qquad
  F=F_{<m}+x^mG_m .
\]
We call $m\in\{1,\ldots,D\}$ a \emph{split} of $F$.  A split is \emph{clean
at} $\alpha$ if $G_m(\alpha)=0$, and \emph{$P$-clean} if it is clean at every
$\alpha_j$.  A polynomial $B\in\mathbb Z[x]$ is \emph{$P$-irreducible} if
$B(0)\ne0$ and no split of $B$ is $P$-clean.

\begin{lemma}[Gaussian carries]\label{lem:gausscarry}
Let $F\in\mathbb Z[x]$ have degree $D$, let $\alpha\in\mathbb Z[i]$ with
$F(\alpha)=0$, $\rho=|\alpha|$, and let $1\le m\le D$ (a split of $F$).  Write
$F=F_{<m}+x^mG_m$ with $F_{<m}=\sum_{i<m}f_ix^i$ and
$G_m=\sum_{i\ge m}f_ix^{i-m}$.  Then
$G_m(\alpha)\in\mathbb Z[i]$ and $F_{<m}(\alpha)=-\alpha^mG_m(\alpha)$.  If the
split $m$ is not clean at $\alpha$, that is, if $G_m(\alpha)\ne0$, then $|F_{<m}(\alpha)|\ge\rho^m$, and if
moreover $\rho>2$, some $i<m$ has
\[
  |f_i|\ \ge\ (\rho/2)^{m-i}.
\]
\end{lemma}

\begin{proof}
$G_m$ has integer coefficients, so $G_m(\alpha)\in\mathbb Z[i]$, and
$0=F(\alpha)=F_{<m}(\alpha)+\alpha^mG_m(\alpha)$.  If $G_m(\alpha)\ne0$ it is a
nonzero Gaussian integer, so $|G_m(\alpha)|\ge1$ and
$|F_{<m}(\alpha)|=\rho^m|G_m(\alpha)|\ge\rho^m$.  If every $i<m$ had
$|f_i|<(\rho/2)^{m-i}$, then
$|F_{<m}(\alpha)|\le\sum_{i<m}|f_i|\rho^i<\rho^m\sum_{i<m}2^{i-m}<\rho^m$, a
contradiction.
\end{proof}

The lemma is a Gaussian analogue of the gap principle for lacunary
polynomials~\cite{lenstra-lacunary}, by which a root of small degree of $F$
is a root of both parts across a sufficiently long run of zero coefficients;
here small coefficients below the split play the role of the gap.  Here it is used at every split,
not only at long gaps.

\begin{lemma}[Blocks]\label{lem:gaussblocks}
In the setting of Section~\ref{sec:gaussian}, every nonzero
$F\in\mathbb Z[x]$ divisible by $P$ can be written
\[
  F=\sum_{s=1}^rx^{e_s}B_s,\qquad e_s+\deg B_s<e_{s+1},
\]
with every $B_s\in\mathbb Z[x]$ divisible by $P$ and $P$-irreducible: $B_s(0)\ne0$
and no split $m\in\{1,\ldots,\deg B_s\}$ of $B_s$ has $G_m(\alpha_j)=0$ for every $j$.
Consequently $\Lambda(F)=\sum_s\Lambda(B_s)$, $|\operatorname{lead}B_s|\ge1$,
and $b_k(F)\ge b_k(B_s)$ for every $s$ and $k$.
\end{lemma}

\begin{proof}
Write $F=x^eF_0$ with $F_0(0)\ne0$; since $P(0)\ne0$, $P\mid F_0$.  Induct
on $\deg F_0$.  If no split of $F_0$ is $P$-clean, $F_0$ itself is the only
block.  Otherwise let $m$ be a $P$-clean split: $G_m(\alpha_j)=0$ for all
$j$, hence $F_{<m}(\alpha_j)=F_0(\alpha_j)-\alpha_j^mG_m(\alpha_j)=0$, and by
reality of the coefficients also $F_{<m}(\bar\alpha_j)=G_m(\bar\alpha_j)=0$.
Since $P$ is monic in $\mathbb Z[x]$ with $2K$ distinct roots, division with
remainder in $\mathbb Z[x]$ shows $P\mid F_{<m}$ and $P\mid G_m$ in
$\mathbb Z[x]$.  Both are nonzero ($F_{<m}(0)=F_0(0)$, and $G_m$ has the leading
coefficient of $F_0$) and of smaller degree than $F_0$; $F_{<m}(0)\ne0$, and
$G_m$ is handled after removing a power of $x$.  Their supports lie in
$[0,m-1]$ and $[m,\deg F_0]$, so the decompositions given by induction
concatenate.  The supports of the $x^{e_s}B_s$ are disjoint and cover that of
$F$, which gives the additivity of $\Lambda$ and the inclusion of the
nonleading magnitudes of each $B_s$ among those of $F$ (the leading
coefficient of $B_s$, $s<r$, is a nonleading coefficient of $F$).
\end{proof}

\begin{theorem}[Degree charging]\label{thm:gausscharge}
In the setting of Section~\ref{sec:gaussian}, assume $\rho_{\min}>2$, and
let $B\in\mathbb Z[x]$ be $P$-irreducible and
divisible by $P$, of degree $d$.  Then
\[
  \sum_{i<d}\log^+|b_i|\ \ge\ d\log\frac{\rho_{\min}}2,
  \qquad\text{so}\qquad
  \Lambda(B)\ \ge\ \log|b_d|+d\log\frac{\rho_{\min}}2 .
\]
Hence every nonzero $F\in\mathbb Z[x]$ divisible by $P$ has
$\Lambda(F)\ge\sum_s(\deg B_s)\log(\rho_{\min}/2)$ over the blocks $B_s$ of
any decomposition of $F$ as in Lemma~\ref{lem:gaussblocks}.
\end{theorem}

\begin{proof}
Every split $m\in\{1,\ldots,d\}$ of $B$ is unclean at some $\alpha_j$, with
$\rho_j\ge\rho_{\min}$, so by Lemma~\ref{lem:gausscarry} there is $i<m$ with
$|b_i|\ge(\rho_{\min}/2)^{m-i}$.  For a fixed $i$, the splits charged to it
this way have distinct $m-i\ge1$ with
$(m-i)\log(\rho_{\min}/2)\le\log|b_i|$, so there are at most
$\log^+|b_i|/\log(\rho_{\min}/2)$ of them.  All $d$ splits are charged, which
is the first inequality.  The rest follows from $\log^+|b_d|\ge\log|b_d|$ and
Lemma~\ref{lem:gaussblocks}.
\end{proof}

Compare Proposition~\ref{prop:noangle}, which lies outside this setting: at
most two roots of $x^{2M}+\rho^{2M}$ above the axis are Gaussian integers
(none unless $\rho^{2M}\in\mathbb Z$).  Indeed, if $\alpha$ and $\alpha'$ are
two of them, $\alpha'/\alpha$ is a root of unity in $\mathbb Q(i)$, so
$\alpha'\in\{\pm\alpha,\pm i\alpha\}$, and at most two of these four points
lie above the axis.  It
pays $\Lambda=2M\log\rho$ for its degree $2M$, a charge of the form the
theorem gives; what the theorem forbids, at Gaussian roots, is a block of
large degree and small mass.

\paragraph{Data availability.}
No datasets were generated or analyzed in this study.

\paragraph{Computer assistance.}
No proof in this paper relies on computation.  The file \texttt{tests/test\_complex.py} in the
repository of~\cite{coefficient-mass} checks the statements, on seeded
random integer multiples and on the extremal examples, each group with a
control showing that a false variant is caught; these checks test the
statements and are not used in any proof.  Polynomial products,
divisibility, coefficient magnitudes, central indices and Rouch\'e
comparisons are computed in exact integer and rational arithmetic; floating
point enters only in evaluating logarithms, in the right-hand sides of mass
bounds and the tropical roots.

\paragraph{Use of AI tools.}
Large language model assistants (Claude, Anthropic) were used in drafting and
revising the text, the proofs and the verification code.  The author has checked the mathematics and is responsible for
the content.

\paragraph{Funding and competing interests.}
The author received no specific funding for this work and has no competing
interests to declare.

\begin{thebibliography}{99}
\small
\bibitem{akian-gaubert-sharify}
M.~Akian, S.~Gaubert and M.~Sharify,
\newblock Log-majorization of the moduli of the eigenvalues of a matrix
polynomial by tropical roots,
\newblock \emph{Linear Algebra and its Applications} 528 (2017), 394--435.

\bibitem{bbbp-ray}
F.~Beaucoup, P.~Borwein, D.~W. Boyd and C.~Pinner,
\newblock Power series with restricted coefficients and a root on a given ray,
\newblock \emph{Mathematics of Computation} 67 (1998), 715--736,
\newblock \href{https://doi.org/10.1090/S0025-5718-98-00960-0}
{doi:10.1090/S0025-5718-98-00960-0}.

\bibitem{coefficient-mass}
B.~Pham,
\newblock Displaced-zero certificates and the coefficient mass of polynomial
multiples,
\newblock preprint, 2026, \url{https://github.com/bangyen/coefficient-mass}.

\bibitem{coefficient-mass-sectors}
B.~Pham,
\newblock Coefficient mass of polynomial multiples in sectors,
\newblock preprint, 2026, \url{https://github.com/bangyen/coefficient-mass}.

\bibitem{dj}
A.~Dubickas and J.~Jankauskas,
\newblock On the reduced height of a polynomial,
\newblock \emph{Publicationes Mathematicae Debrecen} 71 (2007), 325--348,
\newblock \href{https://doi.org/10.5486/PMD.2007.3728}
{doi:10.5486/PMD.2007.3728}.

\bibitem{hayman}
W.~K. Hayman,
\newblock The local growth of power series: a survey of the Wiman--Valiron
method,
\newblock \emph{Canadian Mathematical Bulletin} 17 (1974), 317--358.

\bibitem{lenstra-lacunary}
H.~W. Lenstra, Jr.,
\newblock Finding small degree factors of lacunary polynomials,
\newblock in \emph{Number Theory in Progress}, Vol.~1 (Zakopane, 1997),
de Gruyter, 1999, 267--276.

\bibitem{marden}
M.~Marden,
\newblock \emph{Geometry of Polynomials}, 2nd ed.,
\newblock Mathematical Surveys 3, American Mathematical Society, 1966.

\bibitem{ostrowski}
A.~Ostrowski,
\newblock Recherches sur la m\'ethode de Graeffe et les z\'eros des
polynomes et des s\'eries de Laurent,
\newblock \emph{Acta Mathematica} 72 (1940), 99--155.

\bibitem{schinzel-complex}
A.~Schinzel,
\newblock The reduced length of a polynomial with complex coefficients,
\newblock \emph{Acta Arithmetica} 133 (2008), 73--81,
\newblock \href{https://doi.org/10.4064/aa133-1-5}
{doi:10.4064/aa133-1-5}.
\end{thebibliography}

\end{document}
